\documentclass[11pt]{article}
% Shared preamble for the printable mini lessons (1.4 onward).
%   \documentclass[11pt]{article}
%   % Shared preamble for the printable mini lessons (1.4 onward).
%   \documentclass[11pt]{article}
%   % Shared preamble for the printable mini lessons (1.4 onward).
%   \documentclass[11pt]{article}
%   \input{../preamble}
% Build the key with:  pdflatex -jobname=<name>_KEY "\def\iskey{}\input{<name>}"
% Figures shared with the web page live in figures/Unit_X/<name>.tex and are
% pulled in with \fig{Unit_X/<name>}; the standalone package drops their own
% preamble, so their macros come from figures/preamble.tex, loaded here.
\input{../../figures/preamble}
\usepackage[margin=0.75in]{geometry}
\usepackage{array}
\usepackage{enumitem}
\usepackage{amsmath}
\usepackage[helvet]{sfmath}  % math in the same sans font as the text
\usepackage{standalone}
\pagestyle{empty}
\setlength{\parindent}{0pt}
\setlist[enumerate,1]{leftmargin=1.5em,itemsep=1em}
\setlist[enumerate,2]{label=(\alph*),leftmargin=2em,itemsep=2.2em}

\newif\ifkey
\ifdefined\iskey \keytrue \fi
\colorlet{keyred}{red!75!black}
\newcommand{\ans}[1]{\ifkey\leavevmode{\color{keyred}#1}\fi}
% vertical work space on the student copy, less on the key
\newcommand{\work}[1]{\ifkey\vspace{0.6em}\else\vspace{#1}\fi}
\newcommand{\fig}[1]{\input{../../figures/#1}}
% Title line: \lessontitle{Quartiles, IQR, and Box Plots}
\newcommand{\lessontitle}[1]{{\Large\bfseries #1\ifkey\ \textcolor{keyred}{KEY}\fi}\hfill Name \underline{\hspace{2in}}\par\vspace{0.6em}}
\newcommand{\partner}{\par\vspace{0.4em}Work with a partner. Compare answers when you finish each problem.\par}

% Build the key with:  pdflatex -jobname=<name>_KEY "\def\iskey{}\input{<name>}"
% Figures shared with the web page live in figures/Unit_X/<name>.tex and are
% pulled in with \fig{Unit_X/<name>}; the standalone package drops their own
% preamble, so their macros come from figures/preamble.tex, loaded here.
% Shared preamble for web figures.  Each figure is a standalone .tex that
% \input's this file; scripts/make-figures.sh turns it into an SVG.
\usepackage[T1]{fontenc}
\usepackage{helvet}
\renewcommand{\familydefault}{\sfdefault}
\usepackage{tikz}
\usetikzlibrary{decorations.pathreplacing}
% axis numbers in the text font, not math italic
\pgfkeys{/pgf/number format/.cd,assume math mode=true,1000 sep={}}

% dot plot: \dotplot{min}{max}{value/count,...}{axis label}
\newcommand{\dotplot}[4]{%
\begin{tikzpicture}[x=0.9cm,y=0.33cm]
  \draw[thick] (#1-0.5,0) -- (#2+0.5,0);
  \foreach \t in {#1,...,#2} {\draw (\t,0.25) -- (\t,-0.25) node[below] {\t};}
  \foreach \v/\c in {#3} {\ifnum\c>0 \foreach \k in {1,...,\c} {\fill (\v,\k*1.15) circle (0.16cm);}\fi}
  \node at ({(#1+#2)/2},-2.3) {#4};
\end{tikzpicture}}

% shape sketch histogram (12 bins): \shapehist{heights}{label}
\newcommand{\shapehist}[2]{%
  \foreach \h [count=\i from 0] in {#1} \draw[fill=gray!30] (\i,0) rectangle (\i+1,\h);
  \draw (-0.3,0) -- (12.3,0);
  \node[font=\bfseries] at (6,-2.6) {#2};}
\newcommand{\meanmark}[1]{\fill (#1,0) -- ++(-0.3,-0.75) -- ++(0.6,0) -- cycle;}
\newcommand{\medmark}[2]{\draw[very thick,dashed] (#1,0) -- (#1,#2);}

% number line with minor ticks: \nline{min}{max}{label step}{axis label}
\newcommand{\nline}[4]{%
  \draw[thick] (#1-0.5,0) -- (#2+0.5,0);
  \foreach \t in {#1,...,#2} \draw (\t,0.12) -- (\t,-0.12);
  \pgfmathtruncatemacro{\nxt}{#1+#3}
  \foreach \t in {#1,\nxt,...,#2} \draw (\t,0.2) -- (\t,-0.2) node[below] {\t};
  \node at ({(#1+#2)/2},-0.95) {#4};}
% box plot above the line: \bplot{min}{Q1}{med}{Q3}{max}
\newcommand{\bplot}[5]{%
  \draw[thick] (#1,1.0) -- (#2,1.0);
  \draw[thick] (#4,1.0) -- (#5,1.0);
  \draw[thick] (#1,0.75) -- (#1,1.25);
  \draw[thick] (#5,0.75) -- (#5,1.25);
  \draw[thick,fill=gray!20] (#2,0.55) rectangle (#4,1.45);
  \draw[very thick] (#3,0.55) -- (#3,1.45);}

% ---- added for 1.4 onward ----
% axis on y=0: \axis{min}{max}{label step}{tick step}{axis label}
\newcommand{\axis}[5]{%
  \draw[thick] (#1,0) -- (#2,0);
  \pgfmathsetmacro{\axT}{#1+#4}\foreach \t in {#1,\axT,...,#2} \draw (\t,0.1) -- (\t,-0.1);
  \pgfmathsetmacro{\axL}{#1+#3}\foreach \t in {#1,\axL,...,#2}
    \draw (\t,0.18) -- (\t,-0.18) node[below] {\pgfmathprintnumber{\t}};
  \node at ({(#1+#2)/2},-0.95) {#5};}
% box plot centred at height y: \bplotat{y}{min}{Q1}{med}{Q3}{max}
\newcommand{\bplotat}[6]{%
  \draw[thick] (#2,#1) -- (#3,#1);
  \draw[thick] (#5,#1) -- (#6,#1);
  \draw[thick] (#2,#1-0.22) -- (#2,#1+0.22);
  \draw[thick] (#6,#1-0.22) -- (#6,#1+0.22);
  \draw[thick,fill=gray!20] (#3,#1-0.42) rectangle (#5,#1+0.42);
  \draw[very thick] (#4,#1-0.42) -- (#4,#1+0.42);}
% an outlier dot: \outlier{x}{y}
\newcommand{\outlier}[2]{\fill (#1,#2) circle (2.4pt);}

\usepackage[margin=0.75in]{geometry}
\usepackage{array}
\usepackage{enumitem}
\usepackage{amsmath}
\usepackage[helvet]{sfmath}  % math in the same sans font as the text
\usepackage{standalone}
\pagestyle{empty}
\setlength{\parindent}{0pt}
\setlist[enumerate,1]{leftmargin=1.5em,itemsep=1em}
\setlist[enumerate,2]{label=(\alph*),leftmargin=2em,itemsep=2.2em}

\newif\ifkey
\ifdefined\iskey \keytrue \fi
\colorlet{keyred}{red!75!black}
\newcommand{\ans}[1]{\ifkey\leavevmode{\color{keyred}#1}\fi}
% vertical work space on the student copy, less on the key
\newcommand{\work}[1]{\ifkey\vspace{0.6em}\else\vspace{#1}\fi}
\newcommand{\fig}[1]{\input{../../figures/#1}}
% Title line: \lessontitle{Quartiles, IQR, and Box Plots}
\newcommand{\lessontitle}[1]{{\Large\bfseries #1\ifkey\ \textcolor{keyred}{KEY}\fi}\hfill Name \underline{\hspace{2in}}\par\vspace{0.6em}}
\newcommand{\partner}{\par\vspace{0.4em}Work with a partner. Compare answers when you finish each problem.\par}

% Build the key with:  pdflatex -jobname=<name>_KEY "\def\iskey{}\input{<name>}"
% Figures shared with the web page live in figures/Unit_X/<name>.tex and are
% pulled in with \fig{Unit_X/<name>}; the standalone package drops their own
% preamble, so their macros come from figures/preamble.tex, loaded here.
% Shared preamble for web figures.  Each figure is a standalone .tex that
% \input's this file; scripts/make-figures.sh turns it into an SVG.
\usepackage[T1]{fontenc}
\usepackage{helvet}
\renewcommand{\familydefault}{\sfdefault}
\usepackage{tikz}
\usetikzlibrary{decorations.pathreplacing}
% axis numbers in the text font, not math italic
\pgfkeys{/pgf/number format/.cd,assume math mode=true,1000 sep={}}

% dot plot: \dotplot{min}{max}{value/count,...}{axis label}
\newcommand{\dotplot}[4]{%
\begin{tikzpicture}[x=0.9cm,y=0.33cm]
  \draw[thick] (#1-0.5,0) -- (#2+0.5,0);
  \foreach \t in {#1,...,#2} {\draw (\t,0.25) -- (\t,-0.25) node[below] {\t};}
  \foreach \v/\c in {#3} {\ifnum\c>0 \foreach \k in {1,...,\c} {\fill (\v,\k*1.15) circle (0.16cm);}\fi}
  \node at ({(#1+#2)/2},-2.3) {#4};
\end{tikzpicture}}

% shape sketch histogram (12 bins): \shapehist{heights}{label}
\newcommand{\shapehist}[2]{%
  \foreach \h [count=\i from 0] in {#1} \draw[fill=gray!30] (\i,0) rectangle (\i+1,\h);
  \draw (-0.3,0) -- (12.3,0);
  \node[font=\bfseries] at (6,-2.6) {#2};}
\newcommand{\meanmark}[1]{\fill (#1,0) -- ++(-0.3,-0.75) -- ++(0.6,0) -- cycle;}
\newcommand{\medmark}[2]{\draw[very thick,dashed] (#1,0) -- (#1,#2);}

% number line with minor ticks: \nline{min}{max}{label step}{axis label}
\newcommand{\nline}[4]{%
  \draw[thick] (#1-0.5,0) -- (#2+0.5,0);
  \foreach \t in {#1,...,#2} \draw (\t,0.12) -- (\t,-0.12);
  \pgfmathtruncatemacro{\nxt}{#1+#3}
  \foreach \t in {#1,\nxt,...,#2} \draw (\t,0.2) -- (\t,-0.2) node[below] {\t};
  \node at ({(#1+#2)/2},-0.95) {#4};}
% box plot above the line: \bplot{min}{Q1}{med}{Q3}{max}
\newcommand{\bplot}[5]{%
  \draw[thick] (#1,1.0) -- (#2,1.0);
  \draw[thick] (#4,1.0) -- (#5,1.0);
  \draw[thick] (#1,0.75) -- (#1,1.25);
  \draw[thick] (#5,0.75) -- (#5,1.25);
  \draw[thick,fill=gray!20] (#2,0.55) rectangle (#4,1.45);
  \draw[very thick] (#3,0.55) -- (#3,1.45);}

% ---- added for 1.4 onward ----
% axis on y=0: \axis{min}{max}{label step}{tick step}{axis label}
\newcommand{\axis}[5]{%
  \draw[thick] (#1,0) -- (#2,0);
  \pgfmathsetmacro{\axT}{#1+#4}\foreach \t in {#1,\axT,...,#2} \draw (\t,0.1) -- (\t,-0.1);
  \pgfmathsetmacro{\axL}{#1+#3}\foreach \t in {#1,\axL,...,#2}
    \draw (\t,0.18) -- (\t,-0.18) node[below] {\pgfmathprintnumber{\t}};
  \node at ({(#1+#2)/2},-0.95) {#5};}
% box plot centred at height y: \bplotat{y}{min}{Q1}{med}{Q3}{max}
\newcommand{\bplotat}[6]{%
  \draw[thick] (#2,#1) -- (#3,#1);
  \draw[thick] (#5,#1) -- (#6,#1);
  \draw[thick] (#2,#1-0.22) -- (#2,#1+0.22);
  \draw[thick] (#6,#1-0.22) -- (#6,#1+0.22);
  \draw[thick,fill=gray!20] (#3,#1-0.42) rectangle (#5,#1+0.42);
  \draw[very thick] (#4,#1-0.42) -- (#4,#1+0.42);}
% an outlier dot: \outlier{x}{y}
\newcommand{\outlier}[2]{\fill (#1,#2) circle (2.4pt);}

\usepackage[margin=0.75in]{geometry}
\usepackage{array}
\usepackage{enumitem}
\usepackage{amsmath}
\usepackage[helvet]{sfmath}  % math in the same sans font as the text
\usepackage{standalone}
\pagestyle{empty}
\setlength{\parindent}{0pt}
\setlist[enumerate,1]{leftmargin=1.5em,itemsep=1em}
\setlist[enumerate,2]{label=(\alph*),leftmargin=2em,itemsep=2.2em}

\newif\ifkey
\ifdefined\iskey \keytrue \fi
\colorlet{keyred}{red!75!black}
\newcommand{\ans}[1]{\ifkey\leavevmode{\color{keyred}#1}\fi}
% vertical work space on the student copy, less on the key
\newcommand{\work}[1]{\ifkey\vspace{0.6em}\else\vspace{#1}\fi}
\newcommand{\fig}[1]{\input{../../figures/#1}}
% Title line: \lessontitle{Quartiles, IQR, and Box Plots}
\newcommand{\lessontitle}[1]{{\Large\bfseries #1\ifkey\ \textcolor{keyred}{KEY}\fi}\hfill Name \underline{\hspace{2in}}\par\vspace{0.6em}}
\newcommand{\partner}{\par\vspace{0.4em}Work with a partner. Compare answers when you finish each problem.\par}


\begin{document}

\lessontitle{Choosing Measures and Comparing Distributions}

\textbf{Which measures should you use?} It depends on the shape of the data.

\vspace{0.4em}
\begin{center}
\renewcommand{\arraystretch}{1.4}
\begin{tabular}{|l|c|c|}\hline
\textbf{Shape of the data} & \textbf{Center} & \textbf{Spread} \\ \hline
Symmetric, no outliers & mean & standard deviation \\ \hline
Skewed, or has outliers & median & IQR \\ \hline
\end{tabular}
\end{center}

\textbf{Why?} An outlier pulls the mean and the SD. It does not pull the median and the IQR much.

\vspace{0.4em}
\textbf{Example:} 2, 3, 3, 4, 5, 40

\vspace{0.2em}
\renewcommand{\arraystretch}{1.3}
\begin{tabular}{|l|c|c|c|c|}\hline
 & Mean & SD & Median & IQR \\ \hline
With 40 & 9.5 & 15.0 & 3.5 & 2 \\ \hline
Without 40 & 3.4 & 1.1 & 3 & 2 \\ \hline
\end{tabular}
\qquad
\begin{minipage}{2.6in}
One value moved the mean by 6.1 and the SD by about 14. The median moved 0.5. The IQR did not move.
\end{minipage}

\vspace{1em}
\textbf{Comparing two distributions.} Talk about all four of these. Use numbers and the context.

\vspace{0.3em}
\begin{tabular}{@{}ll}
\textbf{Shape:} & symmetric or skewed?\\
\textbf{Center:} & which group is higher? (median or mean)\\
\textbf{Spread:} & which group is more spread out? (IQR or SD)\\
\textbf{Outliers:} & any values far from the rest?\\
\end{tabular}

\begin{center}\fig{Unit_2/2.2_sleep}\end{center}

\renewcommand{\arraystretch}{1.25}
\begin{tabular}{|>{\raggedright\arraybackslash}p{0.9in}|>{\raggedright\arraybackslash}p{5.7in}|}\hline
Shape & 9th grade is symmetric. 12th grade is skewed right.\\ \hline
Center & The median for 9th grade (8 hours) is greater than for 12th grade (6 hours). 9th graders usually sleep about 2 hours more.\\ \hline
Spread & Both groups have an IQR of 2 hours. The spreads are about the same.\\ \hline
Outliers & One 12th grader slept 12 hours. That is an outlier (the upper fence is 10).\\ \hline
\end{tabular}

\vspace{0.4em}
12th grade has an outlier, so we compare with the median and the IQR.

\partner

% ---------------- PROBLEMS ----------------
\newpage
\begin{enumerate}

\item For each data set, should you use the \textbf{mean and SD} or the \textbf{median and IQR}? Explain.
\begin{enumerate}[itemsep=4em]
  \item The heights of 200 adult women. The histogram is symmetric.

  \ans{Mean and SD. The data is symmetric with no outliers.}
  \item The price of every house sold in a city. The histogram is skewed right.

  \ans{Median and IQR. The data is skewed.}
  \item Quiz scores: 7, 8, 8, 9, 9, 10, 2

  \ans{Median and IQR. 2 is an outlier.}
\end{enumerate}
\work{4em}

\item Look at the example on page 1: 2, 3, 3, 4, 5, 40.
\begin{enumerate}[itemsep=4em]
  \item Which pair describes a typical value better: mean 9.5 and SD 15.0, or median 3.5 and IQR 2? Explain.

  \ans{Median 3.5 and IQR 2. Five of the six values are from 2 to 5. The mean of 9.5 is bigger than all of them except 40.}
  \item A student says, ``The SD is 15, so the values are very spread out.'' Is that a fair description? Explain.

  \ans{No. Five values are packed from 2 to 5. One outlier makes the SD large.}
\end{enumerate}
\work{4em}

\newpage
\item These dot plots show how many books students in two classes read this summer.
\begin{center}\fig{Unit_2/2.2_books}\end{center}
\begin{enumerate}[itemsep=5em]
  \item Describe the shape of each class.

  \ans{Class A is symmetric. Class B is skewed right.}
  \item Class A has median 3 and IQR 2. Class B has median 1 and IQR 2. Compare the centers.

  \ans{The median for Class A (3 books) is greater than for Class B (1 book). Class A usually read 2 more books.}
  \work{1em}
  \item Compare the spreads.

  \ans{Both IQRs are 2 books. The spreads are about the same.}
  \item Does either class have an outlier? Explain.

  \ans{Yes. One student in Class B read 9 books. The upper fence for B is $3 + 1.5 \times 2 = 6$.}
\end{enumerate}

\newpage
\item These box plots show test scores in two class periods.
\begin{center}\fig{Unit_2/2.2_periods}\end{center}
\begin{enumerate}[itemsep=2em]
  \item Compare the centers. Use numbers.

  \ans{The medians are close: Period 2 is 78 and Period 5 is 76. Period 2 is a little higher.}
  \item Compare the spreads. Use numbers.

  \ans{Period 2 IQR $= 85 - 70 = 15$. Period 5 IQR $= 79 - 72 = 7$. Period 2 is much more spread out.}
  \item Which period's scores are more consistent? Explain.

  \ans{Period 5. Its IQR and its range are both smaller.}
\end{enumerate}
\work{2em}

\item Two soccer teams measured the heights of their players. Both histograms are symmetric.

\begin{center}
\renewcommand{\arraystretch}{1.3}
\begin{tabular}{|l|c|c|}\hline
 & Mean & SD \\ \hline
Team X & 70 inches & 2 inches \\ \hline
Team Y & 72 inches & 5 inches \\ \hline
\end{tabular}
\end{center}
\begin{enumerate}[itemsep=2em]
  \item Which team is taller on average? \ans{\quad Team Y, by 2 inches.}
  \item Which team has heights that are more spread out? \ans{\quad Team Y. Its SD is larger.}
  \item Write one sentence that compares the two teams. Use both the center and the spread.

  \ans{Team Y is taller on average (72 vs 70 inches), and its heights are more spread out (SD 5 vs 2 inches).}
\end{enumerate}

\end{enumerate}

\end{document}
